\documentclass[a4paper,11pt]{article}

\usepackage[T1]{fontenc}
\usepackage[utf8]{inputenc}
\usepackage{geometry}
\usepackage{hyperref}
\usepackage{enumitem}
\usepackage{xcolor}
\usepackage{longtable}

\geometry{margin=2.5cm}

\title{
Roadmap Towards the First Practical VML Implementation\\
Discussion Paper for the Technical Working Group for Formal Automotive Systems
}

\author{Working Draft}
\date{\today}

\begin{document}

\maketitle
\begin{center}
\textbf{Origin:} https://twg-fas.org
\end{center}

\begin{abstract}
The VML whitepaper establishes a compelling vision:
natural language is an engineering-time artifact,
while formal machine-interpretable specifications
govern runtime vehicle behavior.

However, a significant gap remains between the conceptual
language definition and a practical interoperable implementation.

This document identifies the most important unresolved questions,
ranks them by implementation relevance,
and proposes a roadmap for future standardization work.

The central question is not whether VML is useful,
but what additional specifications are required so that
two independent organizations can implement VML and obtain
compatible runtime behavior.
\end{abstract}

\tableofcontents

\newpage

%%%%%%%%%%%%%%%%%%%%%%%%%%%%%%%%%%%%%%%%%%%%%%%%%%%%%%%%%%%
\section{Current State of the Initiative}

The current VML draft successfully defines:

\begin{itemize}
\item motivation and architectural role;
\item layered language model;
\item basic grammar;
\item monitor semantics;
\item conformance concepts;
\item examples of behavioral specifications.
\end{itemize}

The current documents answer:

\begin{quote}
``What should VML represent?''
\end{quote}

but do not yet fully answer:

\begin{quote}
``What artifacts must exist,
what data flows through them,
and how is a VML program connected to a real vehicle stack?''
\end{quote}

The following sections identify the highest-priority gaps.

%%%%%%%%%%%%%%%%%%%%%%%%%%%%%%%%%%%%%%%%%%%%%%%%%%%%%%%%%%%
\section{Priority Ranking of Open Topics}

\begin{longtable}{|p{1cm}|p{6cm}|p{6cm}|}
\hline
Priority &
Topic &
Reason \\
\hline

P0 &
Canonical State Model &
No implementation can execute VML without a standardized runtime state representation.\\
\hline

P0 &
Artifact Model &
It is currently unclear what constitutes a VML package or deployment artifact.\\
\hline

P0 &
Signal Binding Specification &
No standard mechanism exists for connecting perception, planning, localization and map data to VML symbols.\\
\hline

P1 &
Evaluation Semantics &
Layer ordering and runtime execution order are not clearly defined.\\
\hline

P1 &
Planner Integration Contract &
Unclear how VML interacts with planners and trajectory generation.\\
\hline

P1 &
Predicate Classification &
No distinction exists between observable facts, derived facts and planner actions.\\
\hline

P2 &
Legal Rule Libraries &
No concept for jurisdiction-specific legal rule sets exists.\\
\hline

P2 &
Ontology Packaging &
Reuse and exchange of ontology definitions remain unspecified.\\
\hline

P3 &
Advanced Probability Model &
Useful but not required for an initial implementation.\\
\hline

P3 &
Certification Profiles &
Important long-term objective but not a blocker for a first prototype.\\
\hline

\end{longtable}

%%%%%%%%%%%%%%%%%%%%%%%%%%%%%%%%%%%%%%%%%%%%%%%%%%%%%%%%%%%
\section{Highest Priority: Canonical State Model}

\subsection{Problem}

The specification states that implementations must map
perception, localization, map and planning data into a
canonical state representation.

However, the structure of that state remains undefined.

Without a common state model there is no interoperability.

Two implementations may use the same VML file
while interpreting symbols differently.

\subsection{Questions to Resolve}

\begin{enumerate}

\item Which primitive types exist?

\item Are coordinates standardized?

\item How are dynamic objects represented?

\item How are lanes represented?

\item How are predicted trajectories represented?

\item How is uncertainty attached to observations?

\item How are timestamps handled?

\item How are multi-rate signals handled?

\end{enumerate}

\subsection{Recommended Deliverable}

VML Runtime State Model v0.1

The state model should become a normative document.

%%%%%%%%%%%%%%%%%%%%%%%%%%%%%%%%%%%%%%%%%%%%%%%%%%%%%%%%%%%
\section{Highest Priority: Artifact Model}

\subsection{Problem}

The specification discusses programs,
parsers,
monitors,
traces,
ontologies
and verification artifacts.

However,
it is unclear what a complete deployable VML package looks like.

\subsection{Open Questions}

Should a VML deployment contain:

\begin{itemize}
\item one file?
\item multiple files?
\item imports?
\item metadata?
\item rule libraries?
\item ontologies?
\item test traces?
\item expected verdict streams?
\end{itemize}

\subsection{Proposed Structure}

\begin{verbatim}
vehicle_mission/
    manifest.yml
    mission.vml
    legal_profile/
    ontology/
    bindings/
    traces/
    verdicts/
\end{verbatim}

A package model should be standardized early.

%%%%%%%%%%%%%%%%%%%%%%%%%%%%%%%%%%%%%%%%%%%%%%%%%%%%%%%%%%%
\section{Signal Binding and Data Mapping}

\subsection{Problem}

The language refers to symbols such as

\begin{verbatim}
Ego.Velocity
DistanceToPedestrian
TrafficLight.State
\end{verbatim}

but their data source is undefined.

\subsection{Questions}

\begin{enumerate}

\item Does VML define signal names?

\item Who owns signal definitions?

\item Are naming conventions mandatory?

\item Can multiple sensor sources provide the same symbol?

\item How is uncertainty attached?

\item How is signal validity represented?

\end{enumerate}

\subsection{Deliverable}

State Binding Specification.

%%%%%%%%%%%%%%%%%%%%%%%%%%%%%%%%%%%%%%%%%%%%%%%%%%%%%%%%%%%
\section{Layer Ordering and Evaluation Semantics}

\subsection{Problem}

The language defines several layers:

\begin{itemize}
\item Logic
\item Temporal
\item Metric
\item Probability
\item Goal
\item Deontic
\item Ontology
\item Utility
\end{itemize}

The specification does not explicitly define whether the ordering
is:

\begin{itemize}
\item pedagogical,
\item conceptual,
\item dependency-based,
\item implementation-based,
\item execution-based.
\end{itemize}

\subsection{Question}

Can two implementations evaluate the same VML program
and guarantee identical results?

\subsection{Recommendation}

Define a normative evaluation pipeline.

Example:

\begin{enumerate}
\item ontology expansion
\item symbol resolution
\item metric evaluation
\item probabilistic evaluation
\item logical evaluation
\item temporal progression
\item deontic activation
\item goal updates
\item utility ranking
\end{enumerate}

%%%%%%%%%%%%%%%%%%%%%%%%%%%%%%%%%%%%%%%%%%%%%%%%%%%%%%%%%%%
\section{Planner Integration Contract}

\subsection{Problem}

The architecture references:

\begin{quote}
VML Evaluator $\rightarrow$
Constraint Filter $\rightarrow$
Trajectory Optimizer
\end{quote}

but the interface remains unspecified.

\subsection{Questions}

\begin{enumerate}

\item Does VML generate planner constraints?

\item Does VML reject trajectories?

\item Does VML assign costs?

\item Does VML activate fallback maneuvers?

\item Can VML request replanning?

\item Can VML force emergency actions?
\end{enumerate}

\subsection{Recommendation}

Define a planner-facing API.

%%%%%%%%%%%%%%%%%%%%%%%%%%%%%%%%%%%%%%%%%%%%%%%%%%%%%%%%%%%
\section{Predicate Classification}

\subsection{Problem}

VML currently treats all predicates similarly.

In practice there are different categories.

\subsection{Proposed Classification}

\begin{description}

\item[Observed Predicate]
Directly measurable.

\item[Derived Predicate]
Computed from state.

\item[Probabilistic Predicate]
Associated with confidence.

\item[Planner Intent Predicate]
Represents planner decisions.

\item[Action Predicate]
Represents requested behavior.

\item[Legal Predicate]
Represents regulatory abstractions.

\end{description}

\subsection{Example}

\begin{verbatim}
TrafficLight.State = RED
\end{verbatim}

Observed Predicate.

\begin{verbatim}
IntersectionClear
\end{verbatim}

Derived Predicate.

\begin{verbatim}
ExecuteMinimalRiskManeuver
\end{verbatim}

Action Predicate.

%%%%%%%%%%%%%%%%%%%%%%%%%%%%%%%%%%%%%%%%%%%%%%%%%%%%%%%%%%%
\section{Jurisdiction and Legal Profiles}

\subsection{Motivation}

A vehicle driving in Germany,
Japan,
California
and Saudi Arabia may follow different rules.

VML therefore requires a mechanism for legal portability.

\subsection{Proposal}

Introduce legal profiles.

\begin{verbatim}
IMPORT LEGAL_PROFILE DE_STVO_2026;
\end{verbatim}

or

\begin{verbatim}
legal_profile:
    jurisdiction: DE
    regulation: STVO
    version: 2026
\end{verbatim}

\subsection{Open Questions}

\begin{itemize}
\item How are legal updates versioned?
\item How is provenance maintained?
\item How are conflicts resolved?
\end{itemize}

%%%%%%%%%%%%%%%%%%%%%%%%%%%%%%%%%%%%%%%%%%%%%%%%%%%%%%%%%%%
\section{End-to-End Example Requirement}

One of the most important missing artifacts is a complete runnable example.

\subsection{Suggested Reference Scenario}

Driving from Ludwigsburg to Stuttgart.

Required artifacts:

\begin{enumerate}

\item Mission specification

\item Route definition

\item State schema

\item Legal rule library

\item Ontology package

\item Example traffic traces

\item Expected monitor outputs

\item Reference implementation

\end{enumerate}

\subsection{Success Criterion}

Two independent implementations
must generate identical verdict streams.

%%%%%%%%%%%%%%%%%%%%%%%%%%%%%%%%%%%%%%%%%%%%%%%%%%%%%%%%%%%
\section{Reference Implementation Roadmap}

\subsection*{Phase 1: Minimal Executable Core}

\begin{itemize}
\item parser
\item AST
\item state schema
\item monitor
\item trace execution
\item verdict generation
\end{itemize}

\subsection*{Phase 2: Automotive Pilot}

\begin{itemize}
\item lane keeping
\item speed limits
\item traffic lights
\item pedestrian protection
\item route missions
\end{itemize}

\subsection*{Phase 3: Ontology and Reuse}

\begin{itemize}
\item inheritance
\item rule libraries
\item legal profiles
\item package management
\end{itemize}

\subsection*{Phase 4: Advanced AI Integration}

\begin{itemize}
\item semantic compilers
\item probabilistic entities
\item risk models
\item planner interaction
\end{itemize}

%%%%%%%%%%%%%%%%%%%%%%%%%%%%%%%%%%%%%%%%%%%%%%%%%%%%%%%%%%%
\section{Recommended Working Group Priorities}

The working group should focus on the following order:

\begin{enumerate}

\item Canonical Runtime State Model

\item Artifact and Package Model

\item Signal Binding Specification

\item Evaluation Semantics

\item Planner Integration Contract

\item Predicate Classification

\item End-to-End Reference Scenario

\item Legal Profiles

\item Ontology Exchange

\item Advanced Probability Extensions

\item Certification Profiles

\end{enumerate}

The first seven topics are likely prerequisites
for meaningful interoperability between independent VML
implementations.

%%%%%%%%%%%%%%%%%%%%%%%%%%%%%%%%%%%%%%%%%%%%%%%%%%%%%%%%%%%
\section{Conclusion}

The current VML initiative has successfully established
a strong conceptual foundation.

The next phase must shift focus from language concepts
toward implementation contracts.

The highest value will likely come from defining:

\begin{enumerate}
\item the canonical runtime state,
\item deployment artifacts,
\item data binding semantics,
\item executable evaluation semantics,
\item planner interfaces.
\end{enumerate}

Once these foundations exist,
VML can evolve from a formal language proposal
into a practical interoperable ecosystem.
The primary objective should therefore be not additional language
features but a precise definition of how a VML specification
interacts with real automotive software systems.

\end{document}